\documentclass{article}
\usepackage[UTF8]{ctex}
\usepackage{amsmath,amssymb,amsfonts}
\usepackage{enumitem}
\usepackage{geometry}
\geometry{a4paper, margin=1in}
% === 错题本排版优化（不参与题目标识） ===
\setlength{\parindent}{0pt}
\setlength{\parskip}{0.35em plus 0.10em minus 0.08em}
\setlength{\emergencystretch}{3em}
\setlist{itemsep=0.25em, topsep=0.40em, parsep=0pt}
\allowdisplaybreaks[2]
% === 错题本排版优化结束 ===
\begin{document}
\section*{武汉理工大学2026年数学分析真题}
\section*{137．武汉理工大学2026年数学分析试卷}

\subsection{题目1}
1．求极限 \(\lim _{n \rightarrow \infty} n^{2}\left(\arctan \frac{2026}{n}-\arctan \frac{2026}{n+1}\right)\) ．

\subsection{题目2}
2．设 \(z=z(x, y)\) 由 \(f(z-x, z-y)=0\) 所确定，其中 \(f(u, v)\) 具有连续的一阶偏导数，且 \(\frac{\partial f}{\partial u}+\frac{\partial f}{\partial v} \neq 0\) ，求 \(\frac{\partial z}{\partial x}+\frac{\partial z}{\partial y}\) 。

\subsection{题目3}
3．已知 \(f(x)=\int_{\sqrt{x}}^{\sqrt{\frac{\pi}{2}}} \frac{1}{1+\tan y^{2}} \mathrm{~d} y\) ，求 \(\int_{0}^{\frac{\pi}{2}} \frac{f(x)}{\sqrt{x}} \mathrm{~d} x\) ．

\subsection{题目4}
4．求第二型曲面积分 \(\iint_{S} \frac{x \mathrm{~d} y \mathrm{~d} z+z \mathrm{~d} x \mathrm{~d} y}{\sqrt{\left(x^{2}+y^{2}+z^{2}\right)^{3}}}\) ，其中 \(S\) 为 \(z=-\sqrt{a^{2}-x^{2}-y^{2}}\) ，方向向上。

\subsection{题目5}
5．求由方程 \(x^{2}+2 x y+2 y^{2}=1\) 确定的隐函数 \(y=y(x)\) 的所有极值，并判断是极大值还是极小值．

\subsection{题目6}
6．求幂级数 \(\sum_{n=1}^{\infty}(-1)^{n-1} n(x-1)^{n}\) 的收敛域与和函数。

\subsection{题目7}
1．已知二元函数

\[
f(x, y)= \begin{cases}\frac{x^{3}+x^{2} y+x y^{2}+2 y^{4}}{x^{2}+y^{2}}, & x^{2}+y^{2} \neq 0 \\ 0, & x^{2}+y^{2}=0\end{cases}
\]

（1）判断 \(f(x, y)\) 在 \((0,0)\) 处的连续性．\\
（2）判断 \(f(x, y)\) 在 \((0,0)\) 处的两个一阶偏导数是否存在？若存在请求其值．\\
（3）判断 \(f(x, y)\) 在 \((0,0)\) 处的可微性，请详细说明理由．

\subsection{题目8}
2．当 \(p>0\) 时，讨论级数 \(\sum_{n=2}^{\infty} \ln \left(1+\frac{(-1)^{n}}{n^{p}}\right)\) 在 \(p\) 为何值时绝对收敛，条件收敛及发散．

\subsection{题目9}
1．证明：函数列 \(f_{n}(x)=\frac{\ln (n x)}{n x^{2}}\) 在 \((1,+\infty)\) 上一致收敛．

\subsection{题目10}
2．证明：\(f(x)\) 在 \((0,1)\) 上一致连续的充分必要条件为对任意 \((0,1)\) 中的柯西列 \(\left\{x_{n}\right\},\left\{f\left(x_{n}\right)\right\}\) 也为柯西列．

\subsection{题目11}
3．证明：数列 \(a_{n}=1+\frac{1}{\sqrt{2}}+\frac{1}{\sqrt{3}}+\cdots+\frac{1}{\sqrt{n}}-2 \sqrt{n}\) 收敛。

\subsection{题目12}
4．解答如下问题：

\begin{itemize}
  \item[（1）] 设 \(k\) 为正整数，证明： \(\int_{(k-1) \pi}^{k \pi} \frac{1}{1+x^{2} \sin ^{2} x} \mathrm{~d} x \geq 2 \int_{0}^{\frac{\pi}{2}} \frac{1}{1+k^{2} \pi^{2} \sin ^{2} x} \mathrm{~d} x\) ．
  \item[（2）] 证明：反常积分 \(\int_{0}^{+\infty} \frac{1}{1+x^{2} \sin ^{2} x} \mathrm{~d} x\) 发散．
\end{itemize}
\end{document}
